\documentclass[11pt,oneside]{article}
\newcommand{\TopicTitle}{Apuntes del Tema 5: Multilayer Perceptron (MLP I y II)}
\usepackage[a4paper,margin=2.6cm,headheight=15pt]{geometry}
\usepackage{fontspec}
\usepackage{polyglossia}
\setdefaultlanguage{spanish}
\setotherlanguage{english}
\setmainfont{Latin Modern Roman}
\setsansfont{Latin Modern Sans}
\setmonofont{Latin Modern Mono}[Scale=MatchLowercase]
\usepackage{amsmath,amssymb,mathtools}
\usepackage{graphicx}
\usepackage{xcolor}
\usepackage{booktabs}
\usepackage{tabularx}
\usepackage{array}
\usepackage{longtable}
\usepackage{enumitem}
\usepackage{listings}
\usepackage{fancyhdr}
\usepackage{titlesec}
\usepackage{microtype}
\usepackage{hyperref}
\usepackage{bookmark}
\usepackage{caption}

\definecolor{ink}{HTML}{243447}
\definecolor{accent}{HTML}{176B87}
\definecolor{codebg}{HTML}{F3F5F7}
\definecolor{codegray}{HTML}{65727E}

\hypersetup{
  colorlinks=true,
  linkcolor=accent,
  urlcolor=accent,
  citecolor=accent,
  pdftitle={\TopicTitle},
  pdfauthor={Foundations of Machine Learning and Deep Learning}
}
\setcounter{secnumdepth}{3}
\setcounter{tocdepth}{3}
\setlist[itemize]{leftmargin=1.7em,itemsep=2pt,topsep=4pt}
\setlist[enumerate]{leftmargin=1.9em,itemsep=2pt,topsep=4pt}
\setlength{\parindent}{0pt}
\setlength{\parskip}{5pt plus 1pt minus 1pt}
\titleformat{\section}{\Large\sffamily\bfseries\color{ink}}{\thesection}{0.8em}{}
\titleformat{\subsection}{\large\sffamily\bfseries\color{accent}}{\thesubsection}{0.8em}{}
\titleformat{\subsubsection}{\normalsize\sffamily\bfseries\color{ink}}{\thesubsubsection}{0.8em}{}
\titlespacing*{\section}{0pt}{2.0em}{0.7em}
\titlespacing*{\subsection}{0pt}{1.4em}{0.4em}
\titlespacing*{\subsubsection}{0pt}{1.0em}{0.3em}
\lstset{
  backgroundcolor=\color{codebg},
  basicstyle=\ttfamily\small,
  breaklines=true,
  breakatwhitespace=false,
  columns=fullflexible,
  frame=single,
  framerule=0.3pt,
  rulecolor=\color{codegray},
  numbers=left,
  numbersep=8pt,
  numberstyle=\scriptsize\color{codegray},
  showstringspaces=false,
  keepspaces=true,
  tabsize=4,
  keywordstyle=\color{accent}\bfseries,
  commentstyle=\color{codegray}\itshape,
  stringstyle=\color{purple}
}
\captionsetup{font=small,labelfont=bf}
\pagestyle{fancy}
\fancyhf{}
\fancyhead[L]{\small\sffamily\color{codegray}Foundations of Machine Learning and Deep Learning}
\fancyhead[R]{\small\sffamily\color{codegray}\nouppercase{\rightmark}}
\fancyfoot[C]{\thepage}
\renewcommand{\headrulewidth}{0.4pt}
\renewcommand{\footrulewidth}{0pt}
\newcommand{\topiccover}[1]{%
  \begin{titlepage}
    \centering
    \vspace*{2.5cm}
    {\Large\sffamily\color{accent}APUNTES UNIVERSITARIOS\par}
    \vspace{1.5cm}
    {\Huge\sffamily\bfseries\color{ink}Foundations of Machine Learning\\[0.25em]
    and Deep Learning\par}
    \vspace{1.8cm}
    {\rule{0.72\textwidth}{1pt}\par}
    \vspace{1.2cm}
    {\LARGE\sffamily\bfseries #1\par}
    \vfill
    {\large\sffamily\color{codegray}\today\par}
  \end{titlepage}
}

\begin{document}
\topiccover{Apuntes del Tema 5: Multilayer Perceptron (MLP I y II)}
\tableofcontents
\clearpage
\markboth{Apuntes del Tema 5: Multilayer Perceptron (MLP I y II)}{Apuntes del Tema 5: Multilayer Perceptron (MLP I y II)}
\section{Apuntes del Tema 5: Multilayer Perceptron (MLP I y II)}

\subsection{1. De Modelos Lineales a No Lineales}

\subsubsection{1.1 Limitaciones de la Transformación Afín}

\begin{itemize}

\item \textbf{Modelos Lineales:} Modelos como la Regresión Lineal o la Regresión Logística calculan transformaciones afines de la forma \(z = x^\top w + b\).

\end{itemize}

\begin{itemize}

\item \textbf{El Problema del Colapso Lineal:} Si se apilan múltiples capas compuestas únicamente por transformaciones afines, la red colapsa matemáticamente en un único modelo lineal equivalente.

\end{itemize}

\begin{itemize}

\item \emph{Demostración:} Dadas dos capas lineales \(H = X W^{(1)} + b^{(1)}\) y \(O = H W^{(2)} + b^{(2)}\):

\end{itemize}

\[O = (X W^{(1)} + b^{(1)}) W^{(2)} + b^{(2)} = X (W^{(1)} W^{(2)}) + (b^{(1)} W^{(2)} + b^{(2)}) = X W + b\]

donde \(W = W^{(1)} W^{(2)}\) y \(b = b^{(1)} W^{(2)} + b^{(2)}\). La capacidad expresiva no aumenta con respecto a un modelo de una sola capa.

\begin{itemize}

\item \textbf{Separabilidad No Lineal (Problema XOR):} Las fronteras de decisión rectilíneas son incapaces de separar clases complejas o disjuntas (como la función lógica XOR o patrones circulares/concéntricos).

\end{itemize}

\subsection{2. Capas Ocultas y Funciones de Activación}

Para dotar al Perceptrón Multicapa (MLP) de la capacidad de aprender relaciones no lineales, se aplica una \textbf{función de activación no lineal} \(\phi(\cdot)\) elemento a elemento tras cada transformación afín en las capas ocultas.

\begin{lstlisting}

Entrada X (n x d) ───► [ Z = X·W(1) + b(1) ] ───► [ H = φ(Z) ] ───► Salida O



\end{lstlisting}

\subsubsection{2.1 Principales Funciones de Activación}

\begin{center}
\begin{tabularx}{\textwidth}{>{\raggedright\arraybackslash}X >{\raggedright\arraybackslash}X >{\raggedright\arraybackslash}X >{\raggedright\arraybackslash}X}
\toprule
Función & Ecuación Matemática \(\phi(z)\) & Rango de Salida & Propiedades y Derivada \(\phi'(z)\) \\
\midrule
\textbf{Sigmoide} & \(\sigma(z) = \frac{1}{1 + e^{-z}}\) & \((0, 1)\) & Interpretación probabilística. \textbf{Problema:} Muestra zonas llanas en los extremos donde la derivada \(\sigma'(z) = \sigma(z)(1 - \sigma(z))\) tiende a 0, provocando desvanecimiento de gradiente (\emph{vanishing gradient}). \\
\bottomrule
\end{tabularx}
\end{center}

 |

| \textbf{Tanh} \emph{(Tangente Hiperbólica)} | \(\tanh(z) = \frac{e^z - e^{-z}}{e^z + e^{-z}}\) | \((-1, 1)\) | Centrada en cero (facilita la optimización). Sufre de desvanecimiento de gradientes para \(\vert{}z\vert{}\) elevados (\(\tanh'(z) = 1 - \tanh^2(z)\)).

 |

| \textbf{ReLU} \emph{(Rectified Linear Unit)} | \(\text{ReLU}(z) = \max(0, z)\) | \([0, \infty)\) | \textbf{Estándar en redes profundas:} Computacionalmente eficiente. Derivada \(\phi'(z) = 1\) para \(z > 0\) (evita el desvanecimiento de gradientes en la región positiva). \textbf{Riesgo:} \emph{Dying ReLU} si \(z \le 0\) de forma permanente. |

\subsection{3. Clasificación Multiclase: Regresión Softmax}

Cuando el target \(y\) pertenece a una de las \(K\) categorías posibles (\(K > 2\)) sin un orden explícito, se emplea la codificación \textbf{One-Hot Encoding} (\(y \in \mathbb{R}^K\)).

\subsubsection{3.1 Operador Softmax}

La capa de salida genera \(K\) valores no acotados denominados \textbf{logits} (\(O = [o_1, o_2, \dots, o_K]\)). Para transformar los logits en una distribución de probabilidad bien definida donde \(\sum_{k=1}^{K} \hat{y}_k = 1\) y \(\hat{y}_k > 0\), se aplica la función \textbf{Softmax}:

\[\hat{y}_k = \text{Softmax}(o_k) = \frac{e^{o_k}}{\sum_{j=1}^{K} e^{o_j}}\]

\begin{itemize}

\item \textbf{Propiedad Conservadora del Orden:} Preserva la ordenación monótona de los logits (\(\arg\max_k \hat{y}_k = \arg\max_k o_k\)).

\end{itemize}

\subsubsection{3.2 Pérdida de Entropía Cruzada Categorical (\emph{Categorical Cross-Entropy})}

Para una muestra con vector One-Hot \(y\) y predicciones \(\hat{y}\):

\[\mathcal{L}(y, \hat{y}) = - \sum_{k=1}^{K} y_k \log(\hat{y}_k)\]

\begin{quote}\textbf{Nota de Implementación (PyTorch):} En la práctica, no se añade manualmente la capa \texttt{Softmax} al final del modelo. Se retorna el tensor de logits directamente y se utiliza \texttt{nn.CrossEntropyLoss()}, la cual combina internamente \texttt{Softmax} y \texttt{Cross-Entropy} de forma numéricamente estable utilizando el truco \emph{log-sum-exp}.\end{quote}

\begin{quote}\end{quote}

\begin{quote}\end{quote}

\subsection{4. Algoritmo de Propagación: Forward Pass y Backpropagation}

\subsubsection{4.1 Forward Pass (Propagación hacia Adelante)}

Para un MLP de \(L\) capas con entrada \(X \in \mathbb{R}^{n \times d}\) (\(H^{(0)} = X\)):

Para \(i = 1, 2, \dots, L-1\):

\[Z^{(i)} = H^{(i-1)} W^{(i)} + b^{(i)}\]

\[H^{(i)} = \phi(Z^{(i)})\]

Salida final (\(L\)):

\[O = H^{(L-1)} W^{(L)} + b^{(L)}\]

\subsubsection{4.2 Backpropagation (Propagación hacia Atrás)}

Consiste en la aplicación sistemática de la \textbf{Regla de la Cadena} multivariante del cálculo desde la función de pérdida \(\mathcal{J}\) hasta las primeras capas de la red para obtener los gradientes \(\frac{\partial \mathcal{J}}{\partial W^{(i)}}\) y \(\frac{\partial \mathcal{J}}{\partial b^{(i)}}\):

\[\frac{\partial \mathcal{J}}{\partial W^{(i)}} = (H^{(i-1)})^\top \cdot \delta^{(i)}, \quad \text{donde } \delta^{(i)} = \frac{\partial \mathcal{J}}{\partial Z^{(i)}}\]

\begin{itemize}

\item \textbf{Gestión de Memoria:} Los tensores intermedios calculados durante el \emph{forward pass} (\(H^{(i)}, Z^{(i)}\)) deben almacenarse en memoria para ser reevaluados durante el \emph{backpropagation}, por lo que entrenar una red neuronal consume considerablemente más memoria que ejecutar inferencia (\texttt{with torch.no\_grad():}).

\end{itemize}

\subsection{5. Propiedades Teóricas e Inicialización de Parámetros}

\subsubsection{5.1 Teorema de Aproximación Universal}

\begin{itemize}

\item \textbf{Definición:} Un Perceptrón Multicapa con tan solo \textbf{una capa oculta} y un número suficiente de neuronas no lineales puede aproximar cualquier función continua en un dominio compacto con un grado de precisión arbitrario \(\epsilon > 0\).

\end{itemize}

\begin{itemize}

\item \textbf{¿Por qué usar redes profundas (más de 1 capa)?} El teorema demuestra la \emph{existencia} teórica de dicha red, pero no especifica la cantidad de neuronas requeridas (que puede ser masiva/exponencial) ni si sus pesos se pueden aprender eficazmente con descenso de gradiente. Añadir profundidad permite representar funciones complejas mediante composiciones jerárquicas con un \textbf{número drásticamente menor de parámetros}.

\end{itemize}

\subsubsection{5.2 Rompimiento de la Simetría e Inicialización}

\begin{itemize}

\item \textbf{Inconveniente de Inicialización en Cero:} Si todos los pesos de una capa oculta se inicializan con la misma constante (p. ej. \(w_{ij} = 0\)), todas las neuronas de esa capa calcularán el mismo valor en el \emph{forward pass} y recibirán exactamente el mismo gradiente en el \emph{backward pass}. El MLP colapsaría al comportamiento de un modelo lineal único.

\end{itemize}

\begin{itemize}

\item \textbf{Estrategias de Inicialización:}

\item \textbf{Xavier / Glorot Initialization:} Diseñada para activaciones sigmoideas o \(\tanh\). Mantiene constante la varianza de las activaciones y gradientes a través de las capas:

\end{itemize}

\[W \sim \mathcal{N}\left(0, \, \frac{2}{n_{\text{in}} + n_{\text{out}}}\right)\]

\begin{itemize}

\item \textbf{He / Kaiming Initialization:} Diseñada para activaciones \textbf{ReLU}, compensando que la mitad de las neuronas se desactivan (\(\le 0\)):

\end{itemize}

\[W \sim \mathcal{N}\left(0, \, \frac{2}{n_{\text{in}}}\right)\]

\subsection{6. Problemas en Redes Profundas: Desvanecimiento y Explosión del Gradiente}

Al encadenar múltiples capas mediante la regla de la cadena, la magnitud del gradiente acumulado en las capas iniciales se calcula como un producto iterativo de matrices de pesos y derivadas de activación:

\[\frac{\partial \mathcal{J}}{\partial W^{(1)}} \propto \prod_{l=2}^{L} W^{(l)} \cdot \phi'(Z^{(l)})\]

\begin{itemize}

\item \textbf{Vanishing Gradients (Desvanecimiento):} Si las derivadas o los valores propios de los pesos son menores a \(1\) (p. ej., con sigmoide/tanh en zonas de saturación), el producto decrece exponencialmente hacia \(0\) conforme la red es más profunda. Las capas iniciales dejan de actualizarse y no aprenden.

\end{itemize}

\begin{itemize}

\item \textbf{Exploding Gradients (Explosión):} Si los valores de los pesos son mayores a \(1\), el producto crece exponencialmente provocando actualizaciones numéricamente inestables (\texttt{NaN} o \texttt{Inf}). Se mitiga aplicando \textbf{Gradient Clipping} (\texttt{nn.utils.clip\_grad\_norm\_}).

\end{itemize}

\subsection{7. Práctica y Aplicación en Código (PyTorch)}

\begin{lstlisting}[language=Python]

import torch

import torch.nn as nn

import torch.optim as optim



# --------------------------------------------------------------------------

# 1. Definición del Modelo MLP Multiclase

# --------------------------------------------------------------------------

class DeepMLP(nn.Module):

    def __init__(self, input_dim=8, hidden_dims=[12, 6], num_classes=3, dropout_rate=0.2):

        super(DeepMLP, self).__init__()

        

        layers = []

        in_dim = input_dim

        

        # Construcción de las capas ocultas con ReLU y Dropout

        for h_dim in hidden_dims:

            linear_layer = nn.Linear(in_dim, h_dim)

            # Inicialización He / Kaiming recomendada para ReLU

            nn.init.kaiming_normal_(linear_layer.weight, nonlinearity='relu')

            nn.init.zeros_(linear_layer.bias)

            

            layers.append(linear_layer)

            layers.append(nn.ReLU())

            layers.append(nn.Dropout(p=dropout_rate))

            in_dim = h_dim

            

        # Capa de salida lineal (entrega logits directos)

        out_layer = nn.Linear(in_dim, num_classes)

        nn.init.xavier_normal_(out_layer.weight)

        nn.init.zeros_(out_layer.bias)

        layers.append(out_layer)

        

        self.network = nn.Sequential(*layers)



    def forward(self, x):

        return self.network(x)



# --------------------------------------------------------------------------

# 2. Bucle de Entrenamiento

# --------------------------------------------------------------------------

def train_mlp_multiclass():

    device = torch.device("cuda" if torch.cuda.is_available() else "cpu")

    

    # Datos sintéticos: 100 muestras, 8 características, 3 clases

    X_dummy = torch.randn(100, 8)

    y_dummy = torch.randint(0, 3, (100,)) # Clases: 0, 1, 2

    

    model = DeepMLP(input_dim=8, hidden_dims=[12, 6], num_classes=3).to(device)

    

    # CrossEntropyLoss combina LogSoftmax + NLLLoss

    criterion = nn.CrossEntropyLoss()

    optimizer = optim.Adam(model.parameters(), lr=0.005)

    

    X_b = X_dummy.to(device)

    y_b = y_dummy.to(device)

    

    model.train()

    for epoch in range(1, 101):

        optimizer.zero_grad()

        

        # 1. Forward Pass (Logits)

        logits = model(X_b)

        loss = criterion(logits, y_b)

        

        # 2. Backward Pass

        loss.backward()

        

        # 3. Gradient Clipping opcional para evitar explosión

        nn.utils.clip_grad_norm_(model.parameters(), max_norm=1.0)

        

        # 4. Actualización de parámetros

        optimizer.step()

        

        if epoch % 25 == 0:

            print(f"Época [{epoch:3d}/100] | Loss: {loss.item():.4f}")



if __name__ == "__main__":

    train_mlp_multiclass()



\end{lstlisting}
\end{document}
